% ====================================================================
% INFORME FINAL - GEOINFORMÁTICA USACH
% Proyecto: TerraMatch
% Segunda Evaluación - Enero 2026
% ====================================================================

\documentclass[11pt,letterpaper]{article}

% ================================
% PAQUETES ESENCIALES
% ================================
\usepackage[utf8]{inputenc}
\usepackage[T1]{fontenc}
\usepackage[spanish,es-tabla]{babel}
\usepackage[left=2.5cm,right=2.5cm,top=2.5cm,bottom=2.5cm]{geometry}

\usepackage{graphicx}
\usepackage{float}
\usepackage{xcolor}
\usepackage{hyperref}
\usepackage{amsmath,amssymb}
\usepackage{booktabs}
\usepackage{caption}
\usepackage{subcaption}
\usepackage{longtable}
\usepackage{array}
\usepackage{multirow}
\usepackage{listings}

\usepackage{fancyhdr}
\usepackage{lastpage}
\usepackage{setspace}

\usepackage[backend=biber,style=apa]{biblatex}
\addbibresource{referencias.bib}

% ================================
% FORMATO GENERAL
% ================================
\onehalfspacing
\setlength{\parskip}{0.6em}
\setlength{\parindent}{0pt}

% ================================
% COLORES INSTITUCIONALES
% ================================
\definecolor{usachblue}{RGB}{0,46,93}
\definecolor{codegreen}{rgb}{0,0.6,0}
\definecolor{codegray}{rgb}{0.5,0.5,0.5}
\definecolor{backcolour}{rgb}{0.95,0.95,0.92}

% ================================
% CONFIGURACIÓN DE CÓDIGO
% ================================
\lstdefinestyle{mystyle}{
    backgroundcolor=\color{backcolour},
    commentstyle=\color{codegreen},
    keywordstyle=\color{magenta},
    numberstyle=\tiny\color{codegray},
    basicstyle=\ttfamily\footnotesize,
    breaklines=true,
    captionpos=b,
    keepspaces=true,
    numbers=left,
    numbersep=5pt,
    tabsize=2,
    language=Python
}
\lstset{style=mystyle}

% ================================
% HYPERREF
% ================================
\hypersetup{
    colorlinks=true,
    linkcolor=usachblue,
    urlcolor=usachblue,
    citecolor=usachblue
}

% ================================
% ENCABEZADO Y PIE
% ================================
\pagestyle{fancy}
\fancyhf{}
\lhead{\textcolor{usachblue}{Geoinformática - TerraMatch}}
\rhead{\textcolor{usachblue}{USACH 2026}}
\cfoot{\thepage\ de \pageref{LastPage}}
\rfoot{\today}

% ====================================================================
% DOCUMENTO
% ====================================================================
\begin{document}

% ====================================================================
% PORTADA
% ====================================================================
\begin{titlepage}
\centering
\vspace*{1cm}

{\Large \textbf{UNIVERSIDAD DE SANTIAGO DE CHILE}}\\
Facultad de Ingeniería\\
Departamento de Ingeniería Informática

\vspace{3cm}

\rule{\textwidth}{0.5mm}

\vspace{0.5cm}
{\LARGE \textbf{TerraMatch}}\\
\vspace{0.3cm}
{\Large Sistema de estimación de satisfacción inmobiliaria basado en información geoespacial}

\vspace{0.5cm}
\rule{\textwidth}{0.5mm}

\vspace{1.5cm}

{\Large \textbf{Segunda Evaluación}}\\
Informe Final

\vspace{1.5cm}

\begin{flushleft}
\textbf{Integrantes del Grupo:}\\
Valentina Del Pilar Barría\\
Felipe Baeza Muñoz\\
Byron Caices Lima\\
Catalina López\\
Jaime Antonio Riquelme Olguín\\
Diego Ignacio Rojas Bustos\\[0.5cm]

\textbf{Profesor:}\\
Francisco Parra O.
\end{flushleft}

\vfill
{\large Enero 2026}
\end{titlepage}

% ====================================================================
% RESUMEN EJECUTIVO
% ====================================================================
\newpage
\section*{Resumen Ejecutivo}
\addcontentsline{toc}{section}{Resumen Ejecutivo}

El mercado inmobiliario chileno enfrenta una crisis de acceso caracterizada por el requisito de un 20\% de pie inicial para créditos hipotecarios, equivalente a aproximadamente 32 veces el ingreso mediano del país. Esta realidad ha generado consumidores cada vez más exigentes que buscan maximizar la satisfacción de su inversión inmobiliaria, demandando herramientas objetivas que permitan evaluar propiedades más allá del precio.

En este contexto, \textbf{TerraMatch} surge como un sistema de apoyo a la decisión que integra información geoespacial para estimar la satisfacción residencial potencial de propiedades en venta. El proyecto se desarrolla sobre cuatro comunas de la Región Metropolitana de Santiago: Estación Central, La Reina, Ñuñoa y Santiago Centro, seleccionadas por su diversidad socioeconómica y heterogeneidad urbana.

La metodología implementada comprende cinco etapas principales: (1) extracción automatizada de 7.702 propiedades mediante web scraping, (2) geocodificación utilizando Google Maps API en EPSG:32719, (3) cálculo de 42 variables espaciales mediante distancias a puntos de interés, (4) diseño de un índice de satisfacción residencial, y (5) entrenamiento de modelos predictivos mediante Machine Learning.

Los resultados demuestran que el modelo LightGBM alcanza un \textbf{R² de 0.8635} con RMSE de 0.3357 y baja autocorrelación espacial de residuos (Moran's I = 0.0695), indicando ausencia de sesgo geográfico. Las variables más importantes incluyen la distancia a centros de salud (50\% de importancia), precio por m² (16\%) y densidad comercial (8\%).

El sistema se ha implementado como una aplicación web funcional con backend en FastAPI y frontend en Nuxt.js, permitiendo a los usuarios explorar propiedades, obtener predicciones de satisfacción y recibir recomendaciones personalizadas según perfiles de usuario (familia, profesional joven, inversionista, adulto mayor).

% ====================================================================
% ÍNDICE
% ====================================================================
\newpage
\tableofcontents
\newpage

% ====================================================================
% INTRODUCCIÓN
% ====================================================================
\section{Introducción}

\subsection{Contexto y Motivación}

El mercado inmobiliario chileno enfrenta crecientes barreras de acceso asociadas al alto costo del pie hipotecario y a restricciones crediticias. El pie hipotecario del 20\% equivale a aproximadamente 32 veces el ingreso mediano del país, lo que excluye a numerosos hogares tanto de créditos bancarios como de subsidios estatales \cite{biobio_ciper_2025}. Entre 2021 y 2023, la demanda de viviendas se contrajo significativamente, con caídas en ventas y mayores desistimientos, observándose en 2025 un stock cercano a 70.000 unidades disponibles.

El estudio IPSOS 2025 evidencia brechas notables de satisfacción entre propietarios (74\%) y arrendatarios (39\%), reforzando la necesidad de herramientas objetivas para comparar alternativas inmobiliarias \cite{ipsos_monitor_2025}. Las decisiones de compra suelen basarse en información fragmentada, centrada en el precio y atributos internos de la vivienda, dejando en segundo plano variables territoriales clave como:

\begin{itemize}
    \item Relación precio/m² y costo total de la propiedad
    \item Seguridad y presencia de servicios de emergencia cercanos
    \item Acceso a transporte público y conectividad urbana
    \item Proximidad a educación, salud y comercio
    \item Acceso a áreas verdes y espacios recreativos
    \item Condiciones de habitabilidad (tamaño, distribución y privacidad)
\end{itemize}

En este contexto surge \textbf{TerraMatch}, un sistema que integra análisis geoespacial y modelamiento predictivo para estimar la satisfacción residencial potencial de propiedades inmobiliarias. La propuesta combina atributos internos de la vivienda con variables geoespaciales del entorno urbano, permitiendo una evaluación objetiva y comparable entre distintas áreas urbanas, reduciendo así la asimetría de información en el mercado.

\subsection{Objetivos}

\subsubsection{Objetivo General}

Desarrollar un sistema de recomendación inmobiliaria que estime la satisfacción residencial potencial, integrando atributos internos de la vivienda con variables geoespaciales del entorno urbano, utilizando técnicas de análisis espacial y \textit{Machine Learning} para apoyar la toma de decisiones de compra en la Región Metropolitana de Santiago.

\subsubsection{Objetivos Específicos}

\begin{enumerate}
    \item Integrar y normalizar 29 datasets geoespaciales en el sistema de referencia EPSG:32719, asegurando control de calidad y consistencia espacial.
    \item Calcular métricas espaciales relevantes (distancias, densidades e índices de accesibilidad) y asociarlas a 7.702 propiedades inmobiliarias geocodificadas.
    \item Diseñar un proxy de satisfacción residencial que combine variables económicas, atributos físicos y características del entorno urbano.
    \item Entrenar y validar un modelo predictivo basado en \textit{Machine Learning}, evaluando su desempeño mediante métricas estadísticas y análisis de autocorrelación espacial.
    \item Implementar un dashboard interactivo que permita explorar los resultados del análisis y facilitar la interpretación de la satisfacción residencial estimada.
\end{enumerate}

\subsection{Preguntas de Investigación}

\subsubsection{Pregunta Principal}
\textbf{¿Qué factores hacen que una vivienda sea satisfactoria para vivir en el contexto urbano de Santiago de Chile?}

\subsubsection{Preguntas Secundarias}
\begin{itemize}
    \item ¿Qué peso tienen los factores espaciales (acceso a servicios, transporte y áreas verdes) respecto al precio y atributos internos de la vivienda?
    \item ¿Existen patrones territoriales y diferencias sistemáticas entre comunas en la satisfacción residencial predicha?
    \item ¿Cómo varían los determinantes de satisfacción según el perfil de usuario (familia, profesional joven, inversionista)?
    \item ¿Es posible predecir la satisfacción residencial con precisión suficiente para apoyar decisiones de compra?
\end{itemize}

\subsubsection{Hipótesis de Trabajo}
\begin{itemize}
    \item \textbf{H1}: Los factores espaciales externos influyen significativamente en la satisfacción residencial, más allá del precio.
    \item \textbf{H2}: Existe autocorrelación espacial positiva en la satisfacción residencial, formando clusters territoriales.
    \item \textbf{H3}: Los determinantes de satisfacción varían según el perfil de usuario considerado.
    \item \textbf{H4}: Incorporar variables geoespaciales mejora la capacidad predictiva frente a modelos basados solo en atributos internos.
\end{itemize}

\subsection{Alcance y Limitaciones del Proyecto}

\textbf{Alcance geográfico:} El estudio cubre cuatro comunas de la Región Metropolitana: La Reina, Ñuñoa, Santiago Centro y Estación Central, representando un gradiente urbano desde zonas residenciales de baja densidad hasta el centro histórico y administrativo.

\textbf{Datos de propiedades:} Se analizaron 7.702 propiedades en venta (5.135 departamentos y 2.567 casas), con precios entre 800 y 73.000 UF y superficies entre 16 y 900 m².

\textbf{Limitaciones principales:}
\begin{itemize}
    \item Cobertura geográfica limitada a 4 comunas de la RM
    \item Ausencia de variables de antigüedad, estado de conservación u orientación
    \item Datos estáticos (2025) sin dinámica temporal del mercado
    \item Sin información de gastos comunes, contribuciones o mantenimiento
    \item Solo propiedades en venta (no arriendos)
\end{itemize}

% ====================================================================
% MARCO TEÓRICO
% ====================================================================
\newpage
\section{Marco Teórico}

El proyecto TerraMatch se fundamenta en tres ejes conceptuales complementarios que permiten abordar el problema de la satisfacción residencial desde una perspectiva integral: (1) accesibilidad urbana y provisión de servicios, (2) economía urbana y modelos hedónicos, y (3) análisis espacial y autocorrelación territorial.

\subsection{Satisfacción Residencial y Calidad de Vida Urbana}

La satisfacción residencial se entiende como la adecuación entre las necesidades del hogar y las condiciones del entorno habitacional. La literatura distingue dimensiones funcionales (acceso a servicios, movilidad, equipamiento) y dimensiones simbólicas (seguridad percibida, identidad barrial, sentido de pertenencia) \cite{marans2015,bonaiuto2015}.

Estudios recientes demuestran que la satisfacción residencial está fuertemente influenciada por factores territoriales que van más allá de las características físicas de la vivienda \cite{park2020}. En particular, la proximidad a servicios urbanos, la calidad del transporte público y el acceso a áreas verdes emergen como determinantes significativos de la calidad de vida percibida.

\subsection{Modelos Hedónicos y Valor de la Vivienda}

Desde la economía urbana, el precio de la vivienda actúa como un agregador de las amenidades y condiciones del entorno \cite{rosen1974,lancaster1966}. Los modelos hedónicos descomponen el valor de la vivienda en atributos internos (superficie, dormitorios, baños) y atributos externos (accesibilidad, calidad ambiental, seguridad), constituyendo una base teórica sólida para interpretar la influencia del entorno en la satisfacción residencial \cite{freeman2003,brueckner2011}.

La aplicación de estos modelos en el contexto latinoamericano ha demostrado que la proximidad a transporte masivo, parques y servicios de salud explica una parte significativa de la variación espacial del valor residencial \cite{alvarez2021,zhang2019}.

\subsection{Accesibilidad Urbana y Servicios}

El concepto de accesibilidad urbana puede modelarse de múltiples formas: distancia mínima, tiempo de viaje, cobertura dentro de radios definidos o densidad de servicios \cite{geurs2004,handy1997}. En ausencia de datos de red vial completos, la distancia euclidiana ofrece un estimador consistente para comparaciones relativas a escala comunal \cite{talen2003,wei2018}.

Los índices compuestos de accesibilidad permiten integrar dimensiones heterogéneas (educación, salud, transporte, comercio) en una escala común, facilitando la interpretación y comparación entre territorios. La definición de radios de influencia se alinea típicamente con umbrales caminables: 300 m (~5 min), 600 m (~10 min) y 1000 m (~15 min).

\subsection{Análisis Espacial y Autocorrelación}

El análisis espacial permite identificar patrones territoriales y clusters de valores similares mediante indicadores como el Índice de Moran \cite{anselin1995}. Este indicador mide la autocorrelación espacial, es decir, el grado en que valores cercanos tienden a ser similares entre sí.

En el contexto de modelos predictivos, verificar la autocorrelación espacial de los residuos es fundamental para validar que el modelo no incorpore sesgos territoriales sistemáticos \cite{fotheringham2002,osullivan2010}. Un valor de Moran's I cercano a cero indica independencia espacial, sugiriendo que el modelo captura adecuadamente la estructura territorial.

\subsection{Machine Learning en Predicción Inmobiliaria}

En la última década, los modelos de aprendizaje automático (Random Forest, Gradient Boosting, LightGBM) han demostrado mejoras sustanciales en precisión frente a enfoques lineales tradicionales, especialmente cuando se integran variables espaciales y socioeconómicas \cite{hoon2020,li2020}.

LightGBM, en particular, ofrece ventajas significativas para datos tabulares con alta dimensionalidad: eficiencia computacional, manejo de variables categóricas y capacidad para capturar relaciones no lineales \cite{medium2025_lightgbm}. Herramientas de explicabilidad como SHAP permiten identificar contribuciones locales y globales de variables, complementando la interpretabilidad de modelos de árboles \cite{lundberg2017}.

% ====================================================================
% ÁREA DE ESTUDIO
% ====================================================================
\newpage
\section{Área de Estudio}

\subsection{Ubicación Geográfica}

La Región Metropolitana de Santiago se ubica en la zona central de Chile, entre 32°55' y 34°19' S y 69°46' y 71°43' O. Es una región sin litoral, con capital en Santiago y un valle intermedio entre la Cordillera de los Andes y la de la Costa. El estudio se centra en cuatro comunas contiguas que representan un gradiente urbano característico de Santiago (Figura \ref{fig:ubicacion}).

\begin{figure}[H]
    \centering
    \includegraphics[width=0.5\textwidth]{figuras/figura1_chile_rm.png}
    \caption{Ubicación de la Región Metropolitana de Santiago en el contexto nacional}
    \label{fig:ubicacion}
\end{figure}

\subsection{Comunas del Estudio}

Las cuatro comunas seleccionadas representan la heterogeneidad del Gran Santiago:

\begin{table}[H]
\centering
\caption{Características de las comunas del área de estudio}
\label{tab:comunas}
\begin{tabular}{@{}llll@{}}
\toprule
\textbf{Comuna} & \textbf{Superficie} & \textbf{Característica Principal} & \textbf{Propiedades} \\
\midrule
Santiago Centro & 22.4 km² & Centro histórico y administrativo & 1.110 \\
Estación Central & 14.1 km² & Alta densificación y conectividad & 2.587 \\
Ñuñoa & 16.9 km² & Residencial mixto consolidado & 2.087 \\
La Reina & 23.4 km² & Baja densidad, mayor ingreso & 1.918 \\
\midrule
\textbf{Total} & 76.8 km² & & \textbf{7.702} \\
\bottomrule
\end{tabular}
\end{table}

\begin{figure}[H]
    \centering
    \includegraphics[width=0.85\textwidth]{figuras/mapa_01_ubicacion_area_estudio.png}
    \caption{Mapa de ubicación del área de estudio: las 4 comunas analizadas en la Región Metropolitana}
    \label{fig:area_estudio}
\end{figure}

\subsection{Justificación de la Selección}

La selección de estas cuatro comunas responde a criterios metodológicos específicos:

\begin{itemize}
    \item \textbf{Diversidad socioeconómica:} Desde estratos altos en La Reina hasta medios-bajos en Estación Central
    \item \textbf{Heterogeneidad urbana:} Zonas de alta densificación vertical versus sectores residenciales tradicionales
    \item \textbf{Contiguidad espacial:} Reduce sesgos por diferencias macro-regionales y facilita comparar condiciones urbanas
    \item \textbf{Disponibilidad de datos:} Alta cantidad de propiedades en venta y buena cobertura de servicios urbanos
\end{itemize}

En términos de movilidad, el área incluye corredores de transporte masivo (Línea 1 y Línea 3 de metro, ejes de buses troncales) que concentran accesibilidad en el centro, mientras que sectores residenciales de menor densidad dependen más del transporte privado y la accesibilidad local.

% ====================================================================
% DATOS Y METODOLOGÍA
% ====================================================================
\newpage
\section{Datos y Metodología}

\subsection{Fuentes de Datos}

Se integraron 29 datasets geoespaciales provenientes de múltiples fuentes, organizados en siete macro-categorías:

\begin{table}[H]
\centering
\caption{Categorías de datos geoespaciales utilizados}
\label{tab:categorias_datos}
\begin{tabular}{@{}lll@{}}
\toprule
\textbf{Categoría} & \textbf{Ejemplos} & \textbf{Fuente} \\
\midrule
Educación & Colegios, jardines, universidades & IDE Chile, MINEDUC \\
Salud & Hospitales, consultorios, farmacias & MINSAL, OSM \\
Transporte & Metro, paraderos de bus, ciclovías & Metro de Santiago, DTPM \\
Seguridad & Comisarías, cuarteles, bomberos & Carabineros, Bomberos \\
Comercio & Centros comerciales, supermercados & OSM, Google Maps \\
Entorno & Parques, plazas, áreas verdes & MINVU, Municipios \\
Propiedades & Casas y departamentos en venta & Portal Inmobiliario \\
\bottomrule
\end{tabular}
\end{table}

\textbf{Datos inmobiliarios:} Se extrajeron 7.702 propiedades en venta mediante web scraping de Portal Inmobiliario, incluyendo precio en UF, superficie útil y total, tipo de vivienda, número de dormitorios y baños, y dirección textual. Las propiedades se geocodificaron utilizando Google Maps API, obteniendo coordenadas en el sistema EPSG:32719 (UTM Zone 19S).

\subsection{Procesamiento de Datos}

El preprocesamiento siguió un pipeline estructurado de control de calidad:

\begin{enumerate}
    \item \textbf{Normalización de CRS:} Todos los datasets reproyectados a EPSG:32719
    \item \textbf{Filtrado espacial:} Selección de elementos dentro del área de estudio con buffer de inclusión
    \item \textbf{Validación geométrica:} Corrección de geometrías inválidas y eliminación de duplicados
    \item \textbf{Estandarización de atributos:} Homogenización de campos y control de valores nulos
    \item \textbf{Geocodificación:} Conversión de direcciones textuales a coordenadas geográficas
\end{enumerate}

\begin{figure}[H]
    \centering
    \includegraphics[width=0.9\textwidth]{figuras/01_distribucion_comunas.png}
    \caption{Distribución de propiedades por comuna y tipo de vivienda}
    \label{fig:distribucion_comunas}
\end{figure}

\subsection{Metodología de Análisis Espacial}

La metodología se organiza en tres fases principales:

\textbf{Fase 1 - Preparación de datos:} Normalización de 29 datasets y consolidación de 7.702 propiedades geocodificadas con sus atributos.

\textbf{Fase 2 - Características espaciales:} Generación de una grilla regular de 3.149 puntos (espaciamiento de 250 m) y cálculo de métricas espaciales:
\begin{itemize}
    \item 21 distancias mínimas a categorías de servicios
    \item 42 densidades de servicios en radios de 300, 600 y 1000 m
    \item 9 índices compuestos de accesibilidad (educación, salud, transporte, seguridad, comercio, entorno)
\end{itemize}

\textbf{Fase 3 - Modelamiento predictivo:} Construcción de un proxy de satisfacción residencial, selección de 42 features y entrenamiento de modelos de Machine Learning con validación cruzada.

\subsubsection{Cálculo de Distancias}

Para cada punto de la grilla se calculó la distancia mínima a cada categoría de servicio utilizando estructuras de datos espaciales (cKDTree), reduciendo la complejidad computacional de O(N×M) a O(N log M):

\begin{equation}
d_{i,k} = \min_{j} \sqrt{(x_{g_i} - x_{s_j})^2 + (y_{g_i} - y_{s_j})^2}
\end{equation}

donde $d_{i,k}$ es la distancia del punto $i$ al servicio más cercano de categoría $k$.

\subsubsection{Cálculo de Densidades}

Las densidades se calcularon contando servicios dentro de radios definidos y normalizando por área:

\begin{equation}
\rho_{i,k,r} = \frac{n_{i,k,r}}{\pi r^2} \times 10^6 \quad \text{[servicios/km²]}
\end{equation}

Los radios seleccionados (300 m, 600 m, 1000 m) corresponden a umbrales caminables que representan 5, 10 y 15 minutos de caminata respectivamente.

\begin{figure}[H]
    \centering
    \includegraphics[width=0.95\textwidth]{figuras/04_distancias_densidades.png}
    \caption{Distribución de distancias y densidades calculadas para las principales categorías de servicios}
    \label{fig:distancias_densidades}
\end{figure}

\subsubsection{Índices de Accesibilidad}

Las distancias y densidades se combinaron en índices compuestos por dimensión, normalizados a una escala de 1 a 10:

\begin{figure}[H]
    \centering
    \includegraphics[width=0.95\textwidth]{figuras/02_indices_principales.png}
    \caption{Índices de accesibilidad por dimensión: educación, salud, transporte, seguridad, comercio y entorno}
    \label{fig:indices_principales}
\end{figure}

Se construyeron además tres índices superiores que integran las dimensiones anteriores:
\begin{itemize}
    \item \textbf{Índice de Vida Urbana:} Combina transporte, comercio y servicios
    \item \textbf{Índice de Calidad de Vida:} Integra salud, educación y entorno
    \item \textbf{Índice de Habitabilidad Global:} Síntesis ponderada de todas las dimensiones
\end{itemize}

\begin{figure}[H]
    \centering
    \includegraphics[width=0.95\textwidth]{figuras/03_indices_superiores.png}
    \caption{Índices superiores: vida urbana, calidad de vida y habitabilidad global}
    \label{fig:indices_superiores}
\end{figure}

\subsection{Diseño del Proxy de Satisfacción}

El proxy de satisfacción residencial se construyó combinando tres componentes:

\begin{enumerate}
    \item \textbf{Valor económico relativo:} Precio por m² normalizado por comuna (invertido, menor precio = mayor satisfacción potencial)
    \item \textbf{Atributos físicos:} Superficie útil, dormitorios por habitante, ratio baños/dormitorios
    \item \textbf{Accesibilidad espacial:} Índices de accesibilidad a servicios urbanos
\end{enumerate}

El proxy se ajusta a cinco perfiles de usuario, modificando los pesos de cada dimensión:

\begin{table}[H]
\centering
\caption{Perfiles de usuario y sus prioridades}
\label{tab:perfiles}
\begin{tabular}{@{}lp{10cm}@{}}
\toprule
\textbf{Perfil} & \textbf{Prioridades} \\
\midrule
Familia con niños & Espacio, educación, seguridad, áreas verdes \\
Profesional joven & Transporte, comercio, precio accesible \\
Inversionista & ROI, transporte, seguridad \\
Adulto mayor & Salud, seguridad, áreas verdes, accesibilidad \\
Balanceado & Equilibrado en todas las dimensiones \\
\bottomrule
\end{tabular}
\end{table}

\subsection{Modelamiento Predictivo}

Se compararon dos algoritmos de Machine Learning:

\begin{itemize}
    \item \textbf{Random Forest:} Modelo baseline con 100 árboles
    \item \textbf{LightGBM:} Gradient Boosting optimizado para datos tabulares
\end{itemize}

La evaluación utilizó partición entrenamiento/prueba (80/20) y validación cruzada de 5 folds, reportando R², RMSE y MAE. Adicionalmente, se verificó la autocorrelación espacial de residuos mediante el Índice de Moran para detectar posibles sesgos territoriales.

\subsection{Herramientas y Librerías}

El pipeline se implementó en Python 3.12 utilizando:

\begin{itemize}
    \item \textbf{Datos espaciales:} GeoPandas, Shapely, pyproj
    \item \textbf{Análisis:} Pandas, NumPy, SciPy
    \item \textbf{Machine Learning:} scikit-learn, LightGBM
    \item \textbf{Visualización:} Matplotlib, Seaborn, Folium
    \item \textbf{Análisis espacial:} PySAL, libpysal
\end{itemize}

% ====================================================================
% APLICACIÓN WEB Y CHATBOT IA
% ====================================================================
\newpage
\section{Aplicación Web y Chatbot de Inteligencia Artificial}

Además del modelo predictivo, TerraMatch implementa una aplicación web completa con un \textbf{chatbot conversacional basado en Inteligencia Artificial} que permite a los usuarios interactuar de manera natural para encontrar propiedades y recibir asesoría financiera personalizada.

\subsection{Arquitectura de la Aplicación}

El sistema se compone de tres capas principales:

\begin{table}[H]
\centering
\caption{Arquitectura técnica de TerraMatch}
\label{tab:arquitectura}
\begin{tabular}{@{}lll@{}}
\toprule
\textbf{Capa} & \textbf{Tecnología} & \textbf{Función} \\
\midrule
Frontend & Nuxt.js 3 + Vue.js & Interfaz de usuario y chatbot \\
Backend & FastAPI + Python & API REST y modelos ML \\
Base de datos & PostgreSQL + PostGIS & Datos espaciales y propiedades \\
IA Conversacional & Gemini 2.5 Flash & Procesamiento de lenguaje natural \\
\bottomrule
\end{tabular}
\end{table}

\subsection{Chatbot Conversacional con Gemini}

El chatbot utiliza el modelo \textbf{Gemini 2.5 Flash} de Google para mantener conversaciones naturales en español chileno. El asistente está diseñado para:

\begin{itemize}
    \item Recopilar preferencias del usuario de forma conversacional (no como formulario)
    \item Validar la capacidad financiera del usuario en tiempo real
    \item Generar búsquedas de propiedades basadas en la conversación
    \item Presentar resultados de manera amigable y personalizada
    \item Ofrecer alternativas cuando el presupuesto no es suficiente
\end{itemize}

\subsubsection{Flujo de Conversación}

El chatbot implementa un flujo estructurado de 6 fases:

\begin{enumerate}
    \item \textbf{Saludo inicial:} Presentación y pregunta sobre tipo de propiedad
    \item \textbf{Recopilación de preferencias:} Comuna, dormitorios, superficie, precio
    \item \textbf{Validación financiera:} Ingreso mensual y capacidad de endeudamiento
    \item \textbf{Búsqueda:} Generación automática de parámetros y consulta al backend
    \item \textbf{Presentación de resultados:} Análisis y recomendación de las mejores opciones
    \item \textbf{Refinamiento:} Ajuste de búsqueda según feedback del usuario
\end{enumerate}

\subsubsection{Ejemplo de Interacción}

\begin{quote}
\textbf{Usuario:} ``Hola, busco un depto en Ñuñoa''

\textbf{Chatbot:} ``¡Hola! Qué buena elección Ñuñoa. Es una comuna muy cotizada. ¿Tienes algún rango de precio en mente? Y cuéntame, ¿cuántos dormitorios necesitas?''

\textbf{Usuario:} ``Unos 3500 UF máximo, 2 dormitorios''

\textbf{Chatbot:} ``Perfecto, 3.500 UF con 2 dormitorios en Ñuñoa es un rango bastante realista. Antes de buscar, ¿podrías contarme aproximadamente cuánto es tu ingreso mensual? Esto me ayuda a verificar que el financiamiento sea viable.''

\textbf{Usuario:} ``Gano como 2 palos''

\textbf{Chatbot:} ``Con \$2.000.000 mensuales, tu cuota máxima sería de unos \$500.000 aprox. Eso te permite un crédito de hasta 3.200 UF aproximadamente. Te recomiendo buscar en el rango de 2.500-3.200 UF para mayor comodidad financiera. ¿Te parece si busco en ese rango?''
\end{quote}

\subsection{Sistema de Validación Financiera}

Una de las funcionalidades más innovadoras del chatbot es su capacidad para \textbf{validar la capacidad financiera del usuario} en tiempo real durante la conversación.

\subsubsection{Parámetros Financieros}

El sistema utiliza reglas bancarias chilenas configurables:

\begin{table}[H]
\centering
\caption{Parámetros de validación financiera}
\label{tab:financiero}
\begin{tabular}{@{}lr@{}}
\toprule
\textbf{Parámetro} & \textbf{Valor} \\
\midrule
Cuota máxima sobre ingreso & 25\% \\
Tasa de interés anual (promedio) & 5.5\% \\
Plazo del crédito hipotecario & 25 años \\
Pie mínimo exigido & 20\% \\
Valor UF referencial & \$37.500 CLP \\
\bottomrule
\end{tabular}
\end{table}

\subsubsection{Cálculo de Capacidad de Endeudamiento}

El sistema calcula automáticamente el precio máximo que el usuario puede pagar utilizando la fórmula de anualidad:

\begin{equation}
\text{Cuota}_{max} = \text{Ingreso} \times 0.25
\end{equation}

\begin{equation}
\text{Monto}_{crédito} = \text{Cuota}_{max} \times \frac{1 - (1 + r)^{-n}}{r}
\end{equation}

\begin{equation}
\text{Precio}_{max} = \frac{\text{Monto}_{crédito}}{1 - \text{Pie}}
\end{equation}

Donde $r$ es la tasa mensual (5.5\%/12) y $n$ es el número de cuotas (25$\times$12 = 300).

\subsubsection{Tabla de Referencia Rápida}

El chatbot utiliza una tabla de referencia para dar respuestas inmediatas:

\begin{table}[H]
\centering
\caption{Capacidad de compra según ingreso mensual}
\label{tab:capacidad}
\begin{tabular}{@{}rrr@{}}
\toprule
\textbf{Ingreso Mensual} & \textbf{Precio Máximo (UF)} & \textbf{Cuota Aprox.} \\
\midrule
\$1.000.000 & 2.000 UF & \$230.000 \\
\$1.500.000 & 3.000 UF & \$350.000 \\
\$2.000.000 & 4.200 UF & \$480.000 \\
\$2.500.000 & 5.500 UF & \$600.000 \\
\$3.000.000 & 6.800 UF & \$720.000 \\
\$4.000.000 & 9.200 UF & \$960.000 \\
\$5.000.000 & 11.500 UF & \$1.200.000 \\
\bottomrule
\end{tabular}
\end{table}

\subsubsection{Penalización por Deudas}

Si el usuario menciona deudas o cupo de tarjetas de crédito, el sistema reduce proporcionalmente su capacidad:

\begin{itemize}
    \item Factor de reducción: 3\% por cada millón de cupo utilizado
    \item Reducción máxima: 30\% de la capacidad total
\end{itemize}

\subsection{Integración con el Backend}

Cuando el chatbot recopila suficiente información, genera automáticamente un \textbf{payload JSON} que se envía al backend para ejecutar la búsqueda:

\begin{lstlisting}[language=Python, caption=Ejemplo de payload generado por el chatbot]
{
  "precio_min": 2500,
  "precio_max": 3200,
  "tipo_inmueble_preferido": "Departamento",
  "comunas_preferidas": ["Nunoa"],
  "dormitorios_min": 2,
  "peso_precio": 0.25,
  "peso_transporte": 0.20,
  "peso_educacion": 0.15,
  "peso_salud": 0.15,
  "peso_areas_verdes": 0.10
}
\end{lstlisting}

\subsection{Presentación de Resultados}

Una vez que el backend retorna las propiedades recomendadas, el chatbot las analiza y presenta de forma amigable:

\begin{itemize}
    \item Resume el número total de propiedades encontradas
    \item Destaca las 3-5 mejores opciones según el score de satisfacción
    \item Menciona puntos fuertes de cada propiedad (cercanía a metro, áreas verdes, etc.)
    \item Ofrece ver más detalles o refinar la búsqueda
    \item Muestra las propiedades en el mapa interactivo
\end{itemize}

\subsection{Tecnologías de Implementación}

El chatbot está implementado en Vue.js 3 con las siguientes características técnicas:

\begin{itemize}
    \item \textbf{SDK de Gemini:} @google/genai (versión 2025)
    \item \textbf{Modelo de IA:} gemini-2.5-flash con temperatura 0.7
    \item \textbf{Contexto máximo:} 2048 tokens de salida
    \item \textbf{Formato de respuesta:} Markdown con bloques de código para payloads
    \item \textbf{Detección de datos financieros:} Expresiones regulares para español chileno
    \item \textbf{Persistencia de sesión:} Historial de conversación mantenido en memoria
\end{itemize}

Esta integración permite que TerraMatch ofrezca una experiencia de usuario completa, combinando la precisión del modelo predictivo con la accesibilidad de una interfaz conversacional natural.

% ====================================================================
% RESULTADOS
% ====================================================================
\newpage
\section{Resultados}

\subsection{Caracterización del Dataset Inmobiliario}

El dataset final comprende 7.702 propiedades distribuidas heterogéneamente entre las cuatro comunas. La composición muestra un predominio de departamentos (66.7\%) sobre casas (33.3\%), reflejando la estructura del mercado inmobiliario en zonas urbanas consolidadas.

\begin{figure}[H]
    \centering
    \includegraphics[width=0.95\textwidth]{figuras/grafico_01_histogramas.png}
    \caption{Distribución de variables principales: precio (UF), superficie útil (m²), dormitorios y baños}
    \label{fig:histogramas}
\end{figure}

Las distribuciones revelan patrones característicos:
\begin{itemize}
    \item \textbf{Precio:} Asimetría positiva con concentración en rangos medios (2.000-5.000 UF)
    \item \textbf{Superficie:} Patrón bimodal asociado a departamentos (40-60 m²) y casas (120-180 m²)
    \item \textbf{Dormitorios:} Predominio de 2-3 dormitorios (72\% del total)
    \item \textbf{Baños:} Mayoría con 1-2 baños (85\% del total)
\end{itemize}

\begin{table}[H]
\centering
\caption{Estadísticas descriptivas del dataset inmobiliario}
\label{tab:estadisticas}
\begin{tabular}{@{}lrrrr@{}}
\toprule
\textbf{Variable} & \textbf{Media} & \textbf{Mediana} & \textbf{Desv. Est.} & \textbf{Rango} \\
\midrule
Precio (UF) & 5.847 & 4.200 & 5.123 & 800 - 73.000 \\
Superficie útil (m²) & 78.4 & 62.0 & 48.7 & 16 - 900 \\
Precio/m² (UF) & 89.2 & 82.5 & 35.4 & 28 - 280 \\
Dormitorios & 2.4 & 2 & 0.9 & 1 - 8 \\
Baños & 1.8 & 2 & 0.7 & 1 - 6 \\
\bottomrule
\end{tabular}
\end{table}

\subsection{Análisis por Comuna}

Las cuatro comunas presentan perfiles inmobiliarios diferenciados:

\begin{figure}[H]
    \centering
    \includegraphics[width=0.95\textwidth]{figuras/grafico_02_analisis_comunas.png}
    \caption{Análisis comparativo de propiedades por comuna: distribución de precios, superficies y características}
    \label{fig:analisis_comunas}
\end{figure}

\begin{figure}[H]
    \centering
    \includegraphics[width=0.95\textwidth]{figuras/10_analisis_por_comuna.png}
    \caption{Análisis detallado de características espaciales por comuna}
    \label{fig:detalle_comunas}
\end{figure}

\textbf{Santiago Centro} presenta la mayor concentración de servicios y el precio por m² más elevado, reflejando su carácter de centro administrativo y comercial. \textbf{Estación Central} destaca por su alta conectividad de transporte (estación intermodal) pero con déficit relativo en áreas verdes. \textbf{Ñuñoa} ofrece un balance entre servicios barriales y calidad residencial, mientras que \textbf{La Reina} presenta mayor heterogeneidad morfológica con predominio de vivienda unifamiliar.

\subsection{Estadísticas de Accesibilidad}

Los índices de accesibilidad revelan diferencias significativas entre comunas y dimensiones:

\begin{table}[H]
\centering
\caption{Índices promedio de accesibilidad por dimensión (escala 1-10)}
\label{tab:indices_accesibilidad}
\begin{tabular}{@{}lrrrr@{}}
\toprule
\textbf{Dimensión} & \textbf{Media} & \textbf{Desv. Est.} & \textbf{Mín} & \textbf{Máx} \\
\midrule
Educación & 5.79 & 1.84 & 1.2 & 9.8 \\
Salud & 4.36 & 2.12 & 0.5 & 9.5 \\
Transporte & 3.60 & 2.66 & 0.1 & 9.9 \\
Entorno & 4.68 & 1.95 & 0.8 & 9.2 \\
Seguridad & 3.74 & 2.08 & 0.3 & 9.1 \\
Comercio & 2.00 & 1.87 & 0.1 & 8.7 \\
\midrule
Habitabilidad Global & 4.03 & 1.52 & 1.1 & 8.4 \\
\bottomrule
\end{tabular}
\end{table}

La dimensión de \textbf{Educación} presenta la mayor accesibilidad promedio (5.79), reflejando la buena cobertura escolar en el área de estudio. En contraste, \textbf{Comercio} muestra el valor más bajo (2.00), indicando concentración de servicios comerciales en zonas específicas.

\begin{figure}[H]
    \centering
    \includegraphics[width=0.85\textwidth]{figuras/05_mapa_habitabilidad.png}
    \caption{Mapa del índice de habitabilidad global calculado a partir de las características espaciales}
    \label{fig:mapa_habitabilidad}
\end{figure}

\subsection{Análisis de Correlaciones}

El análisis de correlaciones entre variables revela relaciones esperadas y consistentes con la teoría:

\begin{figure}[H]
    \centering
    \includegraphics[width=0.90\textwidth]{figuras/grafico_03_correlaciones.png}
    \caption{Matriz de correlaciones entre variables de propiedades y características espaciales}
    \label{fig:correlaciones}
\end{figure}

Hallazgos principales:
\begin{itemize}
    \item Correlación positiva fuerte entre precio y superficie (r = 0.68)
    \item Correlación negativa entre distancia a áreas verdes y precio (r = -0.51)
    \item Los índices de accesibilidad presentan correlaciones moderadas entre sí (r = 0.3-0.5)
    \item El precio por m² correlaciona negativamente con la distancia al centro
\end{itemize}

\subsection{Rendimiento de los Modelos Predictivos}

Se entrenaron y evaluaron dos modelos de Machine Learning para predecir el proxy de satisfacción residencial:

\begin{table}[H]
\centering
\caption{Comparación de rendimiento: Random Forest vs. LightGBM}
\label{tab:comparacion_modelos}
\begin{tabular}{@{}lrrr@{}}
\toprule
\textbf{Modelo} & \textbf{R²} & \textbf{RMSE} & \textbf{MAE} \\
\midrule
Random Forest (baseline) & 0.8431 & 0.3598 & 0.2813 \\
\textbf{LightGBM (final)} & \textbf{0.8635} & \textbf{0.3357} & \textbf{0.2661} \\
\bottomrule
\end{tabular}
\end{table}

El modelo LightGBM superó al baseline Random Forest en todas las métricas, alcanzando un \textbf{R² de 0.8635}, lo que indica que el modelo explica el 86.35\% de la varianza en la satisfacción residencial. La validación cruzada de 5 folds confirmó la estabilidad del modelo (R² = 0.8650 ± 0.0078).

\begin{table}[H]
\centering
\caption{Métricas detalladas del modelo LightGBM final}
\label{tab:metricas_modelo}
\begin{tabular}{@{}lr@{}}
\toprule
\textbf{Métrica} & \textbf{Valor} \\
\midrule
R² (Test) & 0.8635 \\
RMSE (Test) & 0.3357 \\
MAE (Test) & 0.2661 \\
R² Cross-Validation (5-fold) & 0.8650 ± 0.0078 \\
Propiedades de entrenamiento & 6.162 \\
Propiedades de prueba & 1.540 \\
Características predictoras & 42 \\
\bottomrule
\end{tabular}
\end{table}

\begin{figure}[H]
    \centering
    \includegraphics[width=0.85\textwidth]{figuras/grafico_04_dispersion.png}
    \caption{Gráfico de dispersión: satisfacción predicha vs. satisfacción real (R² = 0.8635)}
    \label{fig:dispersion}
\end{figure}

\subsection{Importancia de Variables}

El análisis de importancia de variables revela el peso relativo de cada factor en la predicción:

\begin{figure}[H]
    \centering
    \includegraphics[width=0.95\textwidth]{figuras/grafico_05_importancia_metricas.png}
    \caption{Importancia de las 20 principales variables en el modelo predictivo LightGBM}
    \label{fig:importancia}
\end{figure}

Las variables más importantes son:

\begin{enumerate}
    \item \textbf{Distancia a centros de salud} (50\% de importancia): Factor dominante en la predicción
    \item \textbf{Precio por m²} (16\%): Refleja el valor económico relativo
    \item \textbf{Densidad comercial en 600m} (8\%): Acceso a servicios básicos
    \item \textbf{Distancia a transporte} (7\%): Conectividad urbana
    \item \textbf{Superficie útil} (5\%): Espacio disponible
\end{enumerate}

Este resultado valida la hipótesis H1: las variables espaciales (especialmente salud y transporte) aportan una fracción significativa de la importancia total, más allá del precio.

\subsection{Validación Espacial}

Se evaluó la autocorrelación espacial de los residuos del modelo mediante el Índice de Moran:

\begin{table}[H]
\centering
\caption{Análisis de autocorrelación espacial de residuos}
\label{tab:moran}
\begin{tabular}{@{}lr@{}}
\toprule
\textbf{Indicador} & \textbf{Valor} \\
\midrule
Índice de Moran (I) & 0.0695 \\
p-value & 0.082 \\
Z-score & 1.74 \\
Interpretación & Dependencia espacial baja \\
\bottomrule
\end{tabular}
\end{table}

El valor de Moran's I = 0.0695 (cercano a cero) indica \textbf{ausencia de sesgo geográfico significativo} en las predicciones. El modelo captura adecuadamente la estructura espacial de la satisfacción residencial.

\subsection{Distribución Espacial de la Satisfacción Predicha}

\begin{figure}[H]
    \centering
    \includegraphics[width=0.95\textwidth]{figuras/mapa_03_satisfaccion_predicha.png}
    \caption{Mapa de satisfacción residencial predicha por el modelo LightGBM en las 4 comunas}
    \label{fig:mapa_satisfaccion}
\end{figure}

El mapa revela un patrón espacial claro:
\begin{itemize}
    \item \textbf{Santiago Centro} concentra los valores más altos de satisfacción predicha (~5.5/10 promedio)
    \item \textbf{Estación Central y Ñuñoa} presentan valores intermedios (4.1-4.2/10)
    \item \textbf{La Reina} muestra mayor dispersión con valores más bajos en promedio (3.5/10)
\end{itemize}

Este gradiente centro-periferia es consistente con la distribución de servicios y accesibilidad en el área metropolitana.

% ====================================================================
% DISCUSIÓN
% ====================================================================
\newpage
\section{Discusión}

\subsection{Interpretación de Resultados}

Los resultados obtenidos confirman la hipótesis central del proyecto: \textbf{las variables espaciales influyen significativamente en la satisfacción residencial}, más allá del precio y los atributos físicos de la vivienda. El modelo LightGBM, con un R² de 0.8635, demuestra que es posible predecir la satisfacción residencial con alta precisión utilizando una combinación de factores económicos, físicos y espaciales.

El hallazgo más notable es la \textbf{dominancia de la distancia a centros de salud} (50\% de importancia), un resultado que puede explicarse por varios factores:
\begin{itemize}
    \item La pandemia COVID-19 incrementó la valoración de proximidad a servicios de salud
    \item Los centros de salud actúan como indicadores de urbanización y servicios
    \item La salud es una necesidad básica transversal a todos los perfiles de usuario
\end{itemize}

\subsection{Verificación de Hipótesis}

\begin{table}[H]
\centering
\caption{Verificación de hipótesis de trabajo}
\label{tab:hipotesis}
\begin{tabular}{@{}lll@{}}
\toprule
\textbf{Hipótesis} & \textbf{Estado} & \textbf{Evidencia} \\
\midrule
H1: Variables espaciales significativas & \textbf{Validada} & 65\% de importancia total \\
H2: Autocorrelación espacial positiva & Parcial & Moran's I = 0.07 (débil) \\
H3: Variación por perfil de usuario & \textbf{Validada} & Rankings diferenciados \\
H4: Mejora con variables geoespaciales & \textbf{Validada} & +5\% R² vs. baseline \\
\bottomrule
\end{tabular}
\end{table}

La hipótesis H2 se valida parcialmente: existe autocorrelación espacial pero es débil (Moran's I = 0.0695), lo cual es positivo ya que indica que el modelo captura la mayor parte de la estructura espacial sin dejar patrones residuales significativos.

\subsection{Comparación con Literatura}

Los resultados son consistentes con estudios previos en contextos urbanos similares:

\begin{itemize}
    \item \cite{alvarez2021} reportan que la accesibilidad a servicios explica 40-60\% de la variación en satisfacción residencial en ciudades latinoamericanas
    \item \cite{zhang2019} encuentran que la proximidad a transporte masivo incrementa el valor percibido en 8-15\%
    \item \cite{hoon2020} obtienen R² de 0.75-0.85 con modelos de ML para predicción inmobiliaria
\end{itemize}

Nuestro R² de 0.8635 se sitúa en el rango superior de la literatura, posiblemente debido a la integración de múltiples fuentes de datos espaciales y la granularidad del análisis.

\subsection{Implicancias Prácticas}

Los resultados sugieren prioridades claras para política pública e intervención urbana:

\begin{enumerate}
    \item \textbf{Salud:} Ampliar la cobertura de consultorios y centros de salud fuera del centro
    \item \textbf{Transporte:} Mejorar conectividad en comunas con menor acceso (La Reina)
    \item \textbf{Comercio:} Fomentar desarrollo de servicios barriales en zonas deficitarias
    \item \textbf{Áreas verdes:} Fortalecer espacios públicos en zonas de alta densificación
\end{enumerate}

Para el mercado inmobiliario, el sistema TerraMatch ofrece:
\begin{itemize}
    \item Criterios objetivos para comparar ubicaciones
    \item Recomendaciones personalizadas según perfil de usuario
    \item Transparencia en la evaluación de factores de satisfacción
\end{itemize}

\subsection{Limitaciones del Estudio}

Es importante reconocer las limitaciones del presente trabajo:

\begin{itemize}
    \item \textbf{Cobertura geográfica:} Solo 4 comunas de la RM, dificultando generalización
    \item \textbf{Variables no observadas:} Antigüedad, calidad de construcción, ruido, contaminación
    \item \textbf{Temporalidad:} Datos estáticos de 2025, sin dinámica de mercado
    \item \textbf{Validación subjetiva:} El proxy de satisfacción no ha sido validado con encuestas reales
    \item \textbf{Distancia euclidiana:} No considera tiempos reales de viaje o barreras físicas
\end{itemize}

\subsection{Trabajo Futuro}

Se proponen las siguientes líneas de desarrollo:

\begin{enumerate}
    \item \textbf{Expansión geográfica:} Incorporar más comunas de la RM
    \item \textbf{Variables adicionales:} Integrar datos de delincuencia, contaminación acústica y calidad del aire
    \item \textbf{Validación empírica:} Realizar encuestas de satisfacción real para calibrar el modelo
    \item \textbf{Temporalidad:} Incorporar series de tiempo para capturar dinámica del mercado
    \item \textbf{Red vial:} Calcular tiempos de viaje reales en lugar de distancias euclidianas
\end{enumerate}

% ====================================================================
% CONCLUSIONES
% ====================================================================
\newpage
\section{Conclusiones}

\subsection{Síntesis de Hallazgos}

El proyecto TerraMatch demuestra la viabilidad técnica y el valor práctico de integrar análisis geoespacial con Machine Learning para estimar la satisfacción residencial en el contexto urbano de Santiago de Chile. Los principales hallazgos son:

\begin{enumerate}
    \item \textbf{Las variables espaciales son determinantes:} La distancia a centros de salud, densidad comercial y acceso a transporte explican más del 65\% de la importancia del modelo predictivo.
    
    \item \textbf{Alta capacidad predictiva:} El modelo LightGBM alcanza un R² de 0.8635, superando significativamente enfoques basados solo en precio y atributos físicos.
    
    \item \textbf{Sin sesgo geográfico:} La baja autocorrelación espacial de residuos (Moran's I = 0.0695) indica que el modelo captura adecuadamente la estructura territorial.
    
    \item \textbf{Patrones territoriales claros:} Existe un gradiente centro-periferia en la satisfacción predicha, consistente con la distribución de servicios urbanos.
    
    \item \textbf{Diferenciación por perfil:} Los rankings de propiedades varían según el perfil de usuario, validando el enfoque multi-perfil.
\end{enumerate}

\subsection{Cumplimiento de Objetivos}

\begin{table}[H]
\centering
\caption{Evaluación del cumplimiento de objetivos específicos}
\label{tab:cumplimiento}
\begin{tabular}{@{}lcc@{}}
\toprule
\textbf{Objetivo} & \textbf{Estado} & \textbf{Evidencia} \\
\midrule
Integrar 29 datasets geoespaciales & \checkmark & Sección 4.1 \\
Calcular métricas espaciales & \checkmark & 42 features calculadas \\
Diseñar proxy de satisfacción & \checkmark & 5 perfiles implementados \\
Entrenar modelo predictivo & \checkmark & R² = 0.8635 \\
Implementar dashboard interactivo & \checkmark & Backend + Frontend funcional \\
\bottomrule
\end{tabular}
\end{table}

\subsection{Contribución del Proyecto}

TerraMatch realiza las siguientes contribuciones:

\begin{itemize}
    \item \textbf{Metodológica:} Pipeline reproducible para integración de datos geoespaciales con ML
    \item \textbf{Técnica:} Modelo predictivo de alta precisión para satisfacción residencial
    \item \textbf{Aplicada:} Sistema funcional de recomendación inmobiliaria con API REST
    \item \textbf{Territorial:} Caracterización espacial detallada de 4 comunas de Santiago
\end{itemize}

\subsection{Recomendaciones}

Para usuarios del sistema:
\begin{itemize}
    \item Utilizar el dashboard para comparar propiedades objetivamente
    \item Considerar el perfil de usuario que mejor represente sus necesidades
    \item Complementar las predicciones con visitas presenciales
\end{itemize}

Para tomadores de decisión:
\begin{itemize}
    \item Priorizar inversión en salud y transporte en comunas deficitarias
    \item Fomentar desarrollo comercial barrial en zonas residenciales
    \item Considerar métricas de accesibilidad en planificación urbana
\end{itemize}

\subsection{Reflexión Final}

El proyecto demuestra que la integración de datos geoespaciales con técnicas de Machine Learning ofrece una perspectiva valiosa para entender y predecir la satisfacción residencial. En un mercado inmobiliario cada vez más competitivo y con barreras de acceso crecientes, herramientas como TerraMatch pueden contribuir a reducir la asimetría de información y apoyar decisiones más informadas.

El enfoque desarrollado es replicable a otras ciudades y contextos, abriendo oportunidades para futuras investigaciones y aplicaciones en el ámbito de la geoinformática urbana.

% ====================================================================
% REFERENCIAS
% ====================================================================
\newpage
\printbibliography
\addcontentsline{toc}{section}{Referencias}

% ====================================================================
% ANEXOS
% ====================================================================
\newpage
\appendix
\section{Anexos}

\subsection{Anexo A: Tablas Complementarias}

\begin{table}[H]
\centering
\caption{Detalle de las comunas del área de estudio}
\label{tab:comunas_detalle}
\begin{tabular}{@{}llll@{}}
\toprule
\textbf{Comuna} & \textbf{Coordenadas} & \textbf{Superficie} & \textbf{Altitud} \\
\midrule
Estación Central & 33°27'07''S 70°40'44''O & 14.1 km² & 512 m \\
La Reina & 33°27'00''S 70°33'00''O & 23.4 km² & 680 m \\
Ñuñoa & 33°28'00''S 70°36'00''O & 16.9 km² & 560 m \\
Santiago Centro & 33°27'00''S 70°40'00''O & 22.4 km² & 520 m \\
\bottomrule
\end{tabular}
\end{table}

\begin{table}[H]
\centering
\caption{Fuentes de datos geoespaciales utilizadas}
\label{tab:fuentes}
\begin{tabular}{@{}llll@{}}
\toprule
\textbf{Categoría} & \textbf{Formato} & \textbf{Fecha} & \textbf{Acceso} \\
\midrule
Propiedades inmobiliarias & GeoJSON & 2025 & Web scraping \\
Estaciones de metro & GeoJSON & 2025 & Metro de Santiago \\
Establecimientos educación & GeoJSON & 2024 & MINEDUC \\
Centros de salud & GeoJSON & 2024 & MINSAL \\
Servicios comerciales & GeoJSON & 2025 & OpenStreetMap \\
Áreas verdes & GeoJSON & 2024 & IDE Chile \\
Comisarías y bomberos & GeoJSON & 2024 & Carabineros \\
\bottomrule
\end{tabular}
\end{table}

\subsection{Anexo B: Figuras Adicionales}

\begin{figure}[H]
    \centering
    \includegraphics[width=0.9\textwidth]{figuras/09_estadisticas_descriptivas.png}
    \caption{Panel de estadísticas descriptivas del dataset inmobiliario}
    \label{fig:estadisticas_anexo}
\end{figure}

\begin{figure}[H]
    \centering
    \includegraphics[width=0.9\textwidth]{figuras/07_analisis_correlaciones.png}
    \caption{Análisis detallado de correlaciones entre variables}
    \label{fig:correlaciones_anexo}
\end{figure}

\begin{figure}[H]
    \centering
    \includegraphics[width=0.9\textwidth]{figuras/06_dashboard_resumen.png}
    \caption{Vista del dashboard de TerraMatch}
    \label{fig:dashboard}
\end{figure}

\subsection{Anexo C: Cronograma Ejecutado}

\begin{figure}[H]
    \centering
    \includegraphics[width=0.9\textwidth]{figuras/figura3_carta_gant_sintetica.png}
    \caption{Carta Gantt sintética del proyecto}
    \label{fig:gantt}
\end{figure}

\subsection{Anexo D: Enlaces y Recursos}

\begin{itemize}
    \item \textbf{Repositorio principal:} \url{https://github.com/FelipeBaeza/GeoInformatica}
    \item \textbf{Backend (API):} \url{https://github.com/valecit-a/geo-proyect-backend}
    \item \textbf{Frontend (Dashboard):} \url{https://github.com/Linna-Lpz/geo-proyect-frontend}
    \item \textbf{Datos procesados:} Disponibles en \texttt{autocorrelacion\_espacial/}
\end{itemize}

\end{document}
